% Verified: False
% Refutation notes:
%   The §2.4 (Voice and style modelling) claim overstates the persona feature vector. The draft says the module "computes seven scalar features --- message count, mean word count, formality score (a bounded ratio of formal to casual marker hits), emoji rate, question rate, exclamation rate", but that enumeration lists only six items, and the `VoiceProfile` dataclass in coviber/persona.py (lines 28-38) has exactly six scalar fields: `n_messages`, `avg_words`, `formality`, `emoji_rate`, `question_rate`, `exclaim_rate` (the remaining three fields — `top_openers`, `top_closers`, `vocabulary` — are lists, not scalars). Change "seven" to "six", or add a seventh scalar (e.g. include `top_openers`/`top_closers`/`vocabulary` counts in a different count and re-classify). Secondary nit worth flagging while fixing: the draft describes "the three most-common opener and closer bigrams", but the greeting/closing regexes in persona.py capture variable-length phrases (single tokens like "hi"/"hey"/"thanks" as well as bigrams like "good morning"/"thank you"), and the code takes `most_common(3)` on the raw capture-group strings, not on 2-gram statistics — "opener/closer phrases" would be accurate. Every other technical claim in the section (loader-agnostic core with a single Loader seam, 7 MCP tools listed correctly, ~180-line mcp_server.py — actually 179 lines, per-tool error strings via the try/except in refresh(), COVIBER_CONFIG re-read per call in _settings(), JSONVectorStore default, non-loopback Qdrant stderr warning via _warn_if_remote_qdrant, dedup as content hash, inference-free graph/urgency/persona, bounded linear urgency U(r) ∈ [0,14] at defaults, zero-core-deps with [mcp]/[search]/[qdrant]/[scrape] extras, empty-corpus refusal sentinel in VoiceProfile.system_prompt) is verified against the shipped code.
% Notes to assembler: Cross-section \\ref targets used: sec:arch, sec:workgraph, sec:urgency, sec:persona, sec:record, sec:mcp, sec:vectors, sec:eval, and sec:related:rag/personal/mcp/voice. The assembler must ensure the Architecture / Design section labels these anchors, or replace \\ref{...} with textual references. Pa

\section{Related Work}
\label{sec:related}

CoViber sits at the intersection of four research and engineering threads: retrieval-augmented memory for language models, personal-context replication systems, tool-augmented LLM interfaces, and computational voice modelling. In each we identify the closest neighbors and articulate what CoViber does differently. The recurring axis of difference is that CoViber treats the \emph{source of context} as the only pluggable seam (\S\ref{sec:arch}) and everything downstream --- the cross-source graph (\S\ref{sec:workgraph}), the urgency score (\S\ref{sec:urgency}), the voice profile (\S\ref{sec:persona}) --- as fixed, deterministic, and inference-free.

\subsection{Retrieval-augmented generation and memory frameworks}
\label{sec:related:rag}

The dominant approach to giving LLMs continuity is retrieval-augmented generation (RAG) over a document store, typically driven by dense-embedding similarity~\cite{TODO-rag-lewis-2020,TODO-dense-passage-retrieval}. Application-layer frameworks --- LangChain's memory abstractions~\cite{TODO-langchain-memory}, LlamaIndex's data connectors and index types~\cite{TODO-llamaindex}, and LangGraph's stateful agent graphs~\cite{TODO-langgraph} --- generalize this pattern into pipelines that ingest heterogeneous documents, chunk and embed them, and expose a retriever interface. More recently, dedicated memory services such as mem0~\cite{TODO-mem0} and Zep~\cite{TODO-zep} add long-term episodic memory with LLM-driven summarization and fact extraction on top of a vector index.

CoViber differs from this family along three axes. First, it is \emph{loader-agnostic by construction}: the only seam is the \texttt{Loader} interface (\S\ref{sec:arch}), and every downstream stage --- dedup, work-graph, urgency, semantic search, MCP tools --- operates on the canonical \texttt{Record} schema (\S\ref{sec:record}). The pipeline is a single function that takes a \texttt{Settings} dataclass and dispatches through a registry, so a new source is $\approx$10 lines of code and no core change. LangChain and LlamaIndex are also extensible via connectors, but their downstream primitives (retrievers, agents, tools) still assume a document-centric abstraction; CoViber's downstream is a \emph{work} model (people, projects, channels, tickets), not a document model. Second, CoViber's memory is \emph{not a vector index alone}. Semantic recall is one MCP tool (\texttt{recall}) among seven; the others expose a symbolic work graph (\texttt{who\_is}, \texttt{project\_status}, \texttt{graph\_summary}), a triage queue (\texttt{catch\_me\_up}), a voice profile (\texttt{voice\_profile}), and an ingest hook (\texttt{refresh}). This co-existence of vector and symbolic memory is deliberate: the two answer different questions, and neither subsumes the other. Third, CoViber is \emph{local-first and inference-free at every core stage}. Dedup is a content hash (\S\ref{sec:record}), the graph is deterministic entity extraction with regex primitives (\S\ref{sec:workgraph}), urgency is a bounded linear scoring function (\S\ref{sec:urgency}), and the persona model is closed-form statistics over token counts (\S\ref{sec:persona}). No LLM call is required to update state; the LLM only reads through MCP. This is a strict inversion of the mem0/Zep model, where an LLM is on the write path.

\subsection{Personal-context and continuous-replication systems}
\label{sec:related:personal}

The closest neighbors are systems that model a specific user's context continuously rather than answering document queries. Rewind~\cite{TODO-rewind} records the desktop stream (screen, audio) and indexes it locally for later semantic search; Mem~\cite{TODO-mem-ai} builds a personal knowledge graph from notes and messages with LLM-driven auto-organization; Karpathy's persona notes~\cite{TODO-karpathy-persona} sketch a lightweight, statistically-informed style profile embedded in a system prompt. Enterprise variants include Glean~\cite{TODO-glean} and Microsoft's Copilot memory~\cite{TODO-copilot-memory}, both of which operate as cloud-hosted work-graphs over corporate SSO'd content.

CoViber is architecturally the same shape as Rewind and Mem --- a continuously-updated personal-context store queried through a chat interface --- and its persona component is the direct spiritual descendant of Karpathy's sketch. It differs in three ways. (i) The transport is standardized: CoViber ships an MCP server (\S\ref{sec:mcp}) whose seven tools any conformant client can call, rather than being married to a proprietary desktop app. (ii) The data path is verifiably local: records, work-graph, and vectors live on disk under a single \texttt{data\_dir}, the default vector backend is a plain JSON file (\texttt{JSONVectorStore}, \S\ref{sec:vectors}), and the config parser deliberately refuses shell-style variable expansion so credentials cannot be smuggled into a shared YAML. When a non-loopback Qdrant URL is configured, the store prints a stderr warning at startup so a copy-pasted managed-cluster URL does not silently exfiltrate embeddings. (iii) No LLM is required to run: the core has zero non-stdlib dependencies, and the optional extras (\texttt{[mcp]}, \texttt{[search]}, \texttt{[qdrant]}, \texttt{[scrape]}) each degrade gracefully when absent. In contrast, both Mem and Copilot's memory require a cloud LLM in the write path; Rewind requires its proprietary indexer. The evaluation corpus in this paper is 100\% synthetic (see \S\ref{sec:eval} and the datasheet), which is not merely a legal convenience: it establishes that continuous context replication does not require personal data to be described or benchmarked.

\subsection{MCP and tool-augmented LLMs}
\label{sec:related:mcp}

Tool-augmented language models --- ReAct~\cite{TODO-react-yao-2022}, Toolformer~\cite{TODO-toolformer}, and their descendants --- established that LLMs benefit from structured external calls beyond text generation. The Model Context Protocol (MCP)~\cite{TODO-mcp-spec} generalises this into a client--server contract: any host application (Claude Desktop, Claude Code, custom clients) can discover and invoke tools exposed by any conformant server, with schema and permission negotiated at connection time. Existing MCP servers largely fall into two classes: thin adapters over a single API (Slack, Linear, GitHub) that expose CRUD-style tools, and heavier retrieval servers that wrap a hosted vector index.

CoViber occupies a third position: it is an MCP server whose tools expose \emph{derived} local state --- a work graph, a triage queue, a voice model --- rather than raw API endpoints or a black-box index. The server is a $\approx$180-line FastMCP module (\texttt{coviber/mcp\_server.py}) that re-reads its \texttt{Settings} on every tool call, so a \texttt{COVIBER\_CONFIG} edit takes effect without a restart, and any per-tool failure (unknown loader, missing file, corrupt work-graph JSON) is returned as a readable string rather than propagating a Python traceback across the transport. This design leans on the MCP contract to keep the LLM-facing surface stable while the internal computations evolve: adding a new loader or urgency signal does not change any tool signature. To our knowledge no other public MCP server exposes a persistent, cross-source work graph together with a bounded urgency queue and an inference-free persona through the same session; the closest analogues either publish a raw source (an IMAP MCP server) or a raw retriever (a vector-DB MCP server).

\subsection{Voice and style modelling}
\label{sec:related:voice}

Computational stylometry has a long history in author identification and forensic linguistics~\cite{TODO-stylometry-koppel,TODO-stamatatos-survey}, using features such as function-word frequencies, character n-grams, and average sentence length to fingerprint an author. On the generative side, recent work on style transfer and controllable generation~\cite{TODO-style-transfer-lample,TODO-controllable-generation} conditions an LLM on target-style examples or few-shot exemplars; commercial writing assistants (Grammarly, Lex) fine-tune on user text or use in-context exemplars to imitate a house voice. The intermediate approach --- extracting a compact statistical fingerprint and injecting it as a system-prompt bias --- is the one Karpathy's persona notes~\cite{TODO-karpathy-persona} popularised.

CoViber's \texttt{persona} module (\S\ref{sec:persona}) implements the intermediate approach in a form deliberately restricted to what closed-form statistics can produce. Given the corpus of records the user authored (i.e., records whose \texttt{from\_name} matches the configured identity), it computes seven scalar features --- message count, mean word count, formality score (a bounded ratio of formal to casual marker hits), emoji rate, question rate, exclamation rate --- plus the three most-common opener and closer bigrams and the twenty highest-frequency non-stopword, non-contraction signature words. These are wrapped into a system prompt of the form ``\emph{Write as $\langle$name$\rangle$. Voice: $\langle$register$\rangle$, $\sim$$k$ words per message, $\langle$emoji cue$\rangle$. Typically opens with `$X$' and closes with `$Y$'\ldots'', suitable for injection into any downstream drafting call. The design has three properties absent from LLM-based rewriters. First, it is \emph{deterministic}: identical input produces identical output, byte for byte, so a drafted reply can be regenerated and audited. Second, it is \emph{empty-corpus safe}: if no self-authored records are found, the module returns an explicit refusal-to-fabricate sentinel prompt rather than an implausible ``formal, $\sim$0 words per message'' description that would mislead callers. Third, it is \emph{cheap enough to run on every call}: no training step, no cached model, no GPU. The tradeoff is expressivity --- the profile cannot capture syntactic idiosyncrasies a fine-tuned rewriter could --- but the goal is stylistic guidance to a downstream LLM that already handles fluency, not full imitation.

\smallskip
\noindent In summary, CoViber's contribution is not any single component --- each of \S\ref{sec:related:rag}--\S\ref{sec:related:voice} has stronger prior art on one axis --- but the combination of a loader-agnostic core, a cross-source work model, inference-free scoring and voice modelling, and an MCP-native interface, packaged as a local-first, dependency-light artifact reproducible on synthetic data.

